\chapter*{Kurzfassung}
\addcontentsline{toc}{chapter}{Kurzfassung}
\markboth{Kurzfassung}{Kurzfassung}
\thispagestyle{plain}

An gestapelten mobilen Bauzaunfeldern sind die Stirnflächen der Rahmenrohre zu
detektieren: kleine, untereinander nahezu identische Kreisringe, die dicht
gepackt und wechselseitig verdeckt auftreten.
Für dieses Zielobjekt existieren keine annotierten Datensätze, und die
Handannotation hunderter Kleinobjekte pro Bild ist teuer. Der hier verfolgte
Ansatz umgeht das Problem vollständig: Trainiert wird \emph{ausschließlich} auf
synthetischen Bildern, die aus neun freigestellten Objektausschnitten und sieben
Hintergrundaufnahmen erzeugt werden; validiert und getestet wird auf echten
Aufnahmen.

Damit verschiebt sich der methodische Schwerpunkt von der Netzarchitektur auf die
\textbf{Augmentierungs- und Synthese-Pipeline}: Sie ist hier nicht Beiwerk zu
einem vorhandenen Datensatz, sondern erzeugt die Trainingsverteilung vollständig.
Die Arbeit behandelt daher zunächst die theoretischen Grundlagen der
Datenaugmentierung -- Augmentierung als Kodierung von Invarianzen, als
Rauschregularisierung und als Mittel gegen \enquote{Shortcut Learning} -- und
entwickelt daraus eine Begründungspflicht: Jeder Parameter stützt sich auf genau
eine von drei Quellen -- eine an den Realdaten \emph{gemessene} Eigenschaft, eine
\emph{physikalische Randbedingung} des Aufnahmegegenstands oder ein in der
\emph{Literatur} belegtes Argument; liegt keine davon vor, wird der Wert
ausdrücklich als \emph{Annahme} gekennzeichnet. Ziel ist, jede Abkürzung zu
schließen, über die das Netz sonst \enquote{hier wurde etwas eingeklebt} anstelle
von \enquote{hier ist ein Rohrende} lernen könnte. Die resultierende Pipeline
umfasst 36 dokumentierte Transformationen mit 40 Parameterwerten (6000 Bilder,
173\,100 gerenderte Instanzen, 136\,384 gelabelte Objekte): 9 sind gemessen,
8 literaturgestützt, 3 physikalisch begründet, 9 gemischt und \textbf{11
ausdrücklich als Annahme gekennzeichnet}.

Auf zwei unberührten realen Aufnahmen erreicht das Verfahren \mAPfifty{} von
\textbf{0,886} und \textbf{0,956} bei einem besten $F_1$ von 0,862 und 0,971;
auf \code{test} sitzt dabei jede Detektion auf einem annotierten Objekt
(Precision 1,000). Die Bewertung liefert darüber hinaus vier Befunde. Erstens ist
der Sim-to-Real-Abstand messbar und beträgt 9,6 bzw.\ 2,6\PP{} \mAPfifty{}
(0,982 synthetisch gegen 0,886/0,956 real) -- klein genug, dass das Verfahren
funktioniert, und weit entfernt von dem Muster, das reines Artefaktlernen
erzeugen würde. Zweitens dominiert nicht eine
Augmentierungsentscheidung, sondern die Inferenzauflösung: \mAPfifty{} steigt von
0,435 bei 640\px{} auf 0,904 bei 1600\px. Drittens verbessert stärkere
Verdeckung in den Trainingsdaten das Ergebnis in vier von fünf Metrikspalten
\emph{monoton} -- obwohl der stärkste Arm 33\,\% weniger Labels enthält.
Viertens deckt die Fehleranalyse einen systematischen Größenversatz von 1,2 bis
2,3\px{} zwischen Generatorlabel und Handannotation auf. Er erklärt den auffällig
niedrigen \mAPseven{}-Wert auf dem Testbild -- eine Sensitivitätsrechnung hebt
ihn von 0,326 auf 0,752, ohne dass am Modell etwas geändert wird -- und lässt
\mAPfifty{} und das $F_1$ unberührt. Bei programmatisch erzeugten
Labels wird die Labelkonvention damit selbst zu einer messbaren Größe.
